%--------------------------------------------------------------------------------
%
%
%--------------------------------------------------------------------------------
\pagebreak
\section*{XOR Boolean Operation}
\addcontentsline{toc}{section}{XOR Boolean Operation}

\instructionentry{Syntax}{
  \texttt{XOR[.N] RegR, RegB, RegA} \\
  \texttt{XOR[.N] RegR, RegB, SignedImmed15}
  
  \texttt{OR[.N/D/W/H/B] RegR, SignedImmed13( RegB )} \\
  \texttt{OR[.N/D/W/H/B] RegR, RegX ( RegB )}
}

\instructionentry{Format}{

	The XOR instruction uses the register, the immediate, the indexed and the register indexed instruction formats. \\

   	\begin{flushleft}
    \begin{tikzpicture}[scale=0.7, transform shape]
        
        \draw[help lines, gray!70, dashed] (0,0) grid(16,1.5);
        
        \node[  tsRectangle, 
                minimum width=16cm,
                minimum height=1cm,
                text centered,
                fill=white!50] at (8,1) { };
                
     	\draw[-] (1,0.5) -- (1, 1.5);
     	\draw[-] (3,0.5) -- (3, 1.5);
     	\draw[-] (5,0.5) -- (5, 1.5);
     	\draw[-] (5.5,0.5) -- (5.5, 1.5);
     	\draw[-] (6,0.5) -- (6, 1.5);
     	\draw[-] (6.5,0.5) -- (6.5, 1.5);
     	\draw[-] (8.5,0.5) -- (8.5, 1.5);
     	\draw[-] (9.5,0.5) -- (9.5, 1.5);
     	\draw[-] (11.5,0.5) -- (11.5, 1.5);
     	
     	\node at (0.5,0.25) {2};
    	\node at (2,0.25) {4};
    	\node at (4,0.25) {4};
    	\node at (5.75,0.25) {3};
    	\node at (7.5,0.25) {4};
    	\node at (9,0.25) {2};
    	\node at (10.5,0.25) {4};
   	 	\node at (13.75,0.25) {9};
   	 	
   	 	\node at (0.5,1) {\small 0};
		\node at (2,1) {\small \OpCodeOR};
		\node at (4,1) {\small Reg R};
		\node at (5.25,1) {\small N};
		\node at (5.75,1) {\small 0};
		\node at (6.25,1) {\small 0};
		\node at (7.5,1) {\small Reg B};
		\node at (9,1) {\small 0};
		\node at (10.5,1) {\small Reg A};
		\node at (13.75,1) {\small 0 };
   	 	
       \end{tikzpicture}
	\end{flushleft}
 	\medskip
 	\begin{flushleft}
    	\begin{tikzpicture}[scale=0.7, transform shape]
        
        % \draw[help lines, gray!70, dashed] (0,0) grid(16,1.5);
        
        \node[  tsRectangle, 
                minimum width=16cm,
                minimum height=1cm,
                text centered,
                fill=white!50] (blocks) at (8,1) { };
                
     	\draw[-] (1,0.5) -- (1, 1.5);
     	\draw[-] (3,0.5) -- (3, 1.5);
     	\draw[-] (5,0.5) -- (5, 1.5);
     	\draw[-] (5.5,0.5) -- (5.5, 1.5);
     	\draw[-] (6,0.5) -- (6, 1.5);
     	\draw[-] (6.5,0.5) -- (6.5, 1.5);
     	\draw[-] (8.5,0.5) -- (8.5, 1.5);
               	
     	\node at (0.5,0.25) {2};
    	\node at (2,0.25) {4};
    	\node at (4,0.25) {4};
    	\node at (5.75,0.25) {3};
    	\node at (7.5,0.25) {4};
   	 	\node at (11.5,0.25) {15};
   	 	
   	 	\node at (0.5,1) {\small 0};
		\node at (2,1) {\small \OpCodeOR};
		\node at (4,1) {\small Reg R};
		\node at (5.25,1) {\small N};
		\node at (5.75,1) {\small 0};
		\node at (6.25,1) {\small 1};
		\node at (7.5,1) {\small Reg B};
		\node at (11.5,1) {\small S-IMM-15};
   	 	
       \end{tikzpicture}
	\end{flushleft}
	\medskip
	\begin{flushleft}
        \begin{tikzpicture}[scale=0.7, transform shape]
        
        % \draw[help lines, gray!70, dashed] (0,0) grid(16,1.5);
        
        \node[  tsRectangle, 
                minimum width=16cm,
                minimum height=1cm,
                text centered,
                fill=white!50] (blocks) at (8,1) { };
                
     	\draw[-] (1,0.5) -- (1, 1.5);
     	\draw[-] (3,0.5) -- (3, 1.5);
     	\draw[-] (5,0.5) -- (5, 1.5);
     	\draw[-] (5.5,0.5) -- (5.5, 1.5);
     	\draw[-] (6,0.5) -- (6, 1.5);
     	\draw[-] (6.5,0.5) -- (6.5, 1.5);
     	\draw[-] (8.5,0.5) -- (8.5, 1.5);
     	\draw[-] (9.5,0.5) -- (9.5, 1.5);
     	
     	\node at (0.5,0.25) {2};
    	\node at (2,0.25) {4};
    	\node at (4,0.25) {4};
    	\node at (5.75,0.25) {3};
    	\node at (7.5,0.25) {4};
    	\node at (9,0.25) {2};
     	\node at (12.25,0.25) {13};
   	 	
   	 	\node at (0.5,1) {\small 1};
		\node at (2,1) {\small \OpCodeOR};
		\node at (4,1) {\small Reg R};
		\node at (5.25,1) {\small N};
		\node at (5.75,1) {\small 0};
		\node at (6.25,1) {\small 0};	
	    \node at (7.5,1) {\small Reg B};
		\node at (9,1) {\small dw};
		\node at (12.25,1) {\small S-IMM-13};
   	 	
       \end{tikzpicture}
	\end{flushleft}
 	\medskip
	\begin{flushleft}
        \begin{tikzpicture}[scale=0.7, transform shape]
        
        % \draw[help lines, gray!70, dashed] (0,0) grid(16,1.5);
        
        \node[  tsRectangle, 
                minimum width=16cm,
                minimum height=1cm,
                text width=3cm,
                text centered,
                fill=white!50] (blocks) at (8,1) { };
                
     	\draw[-] (1,0.5) -- (1, 1.5);
     	\draw[-] (3,0.5) -- (3, 1.5);
     	\draw[-] (5,0.5) -- (5, 1.5);
     	\draw[-] (5.5,0.5) -- (5.5, 1.5);
     	\draw[-] (6,0.5) -- (6, 1.5);
     	\draw[-] (6.5,0.5) -- (6.5, 1.5);
     	\draw[-] (8.5,0.5) -- (8.5, 1.5);
     	\draw[-] (9.5,0.5) -- (9.5, 1.5);
     	\draw[-] (11.5,0.5) -- (11.5, 1.5);
     	
     	\node at (0.5,0.25) {2};
    	\node at (2,0.25) {4};
    	\node at (4,0.25) {4};
    	\node at (5.75,0.25) {3};
    	\node at (7.5,0.25) {4};
    	\node at (9,0.25) {2};
    	\node at (10.5,0.25) {4};
   	 	\node at (13.75,0.25) {9};
   	 	
   	 	\node at (0.5,1) {\small 1};
		\node at (2,1) {\small \OpCodeOR};
		\node at (4,1) {\small Reg R};
		\node at (5.25,1) {\small N};
		\node at (5.75,1) {\small 0};
		\node at (6.25,1) {\small 1};
		\node at (7.5,1) {\small Reg B};
		\node at (9,1) {\small dw};
		\node at (10.5,1) {\small Reg X};
		\node at (13.75,1) {\small 0};
   	 	
       \end{tikzpicture}
	\end{flushleft}
}

\instructionentry{Description}{

The \texttt{XOR} instruction performs the binary XOR operation on two 64-bit operands and stores the result to the target register \texttt{RegR}. The result can be negated by setting the "N" bit. The instruction supports the immediate, the register and the two memory reference instruction formats. The boolean operation itself does not cause any traps.The indexed formats may cause a memory reference associated trap.
}

\instructionentry{Operation}{

  	\texttt{RegR <- RegB \^{} RegA} {(register format)}\\
  	\smallskip
  	\texttt{RegR <- RegB \^{} immOperand( )} {(immediate format)}\\
   	\smallskip
  	\texttt{RegR <- RegR \^{} memOperand( )} {(indexed formats)}\\
  	\texttt{if ( instr.[N] RegR <- not( RegR ); }
}

\instructionentry{Exceptions}{

	Data Alignment Trap\\
	Memory Access Violation Trap\\
	Memory Protection Violation Trap\\
	Data TLB Miss Trap
}

\instructionentry{Notes}{

	None. 
}




